\chapter{The Interaction Family}
\label{ch:gaussian}

The aggregation, Kron, and operator results of this supplement act on
precisions of one shape: a matrix-weighted graph Laplacian plus positive
semidefinite self terms. This chapter asks when the belief-sector precision of
\Cref{ch:generative-gaussian} has that shape. The answer is that the shape is
a restriction, not a consequence of linear-Gaussian directedness: an
orientation-reversing gain, and for odd \(K\) a v-structure, leave it in every
chart reached by orientation-preserving frames
(\Cref{prop:gauss-law-level-witnesses}), while on forests every
orientation-preserving invertible gain is absorbed
(\Cref{prop:gauss-gain-absorption}); the condition for membership is
\Cref{prop:gauss-global-interaction-condition}, and the restriction is adopted
as \Cref{hyp:gauss-global-interaction}. It is used by hard aggregation
(\Cref{sec:cg-gaussian-aggregation}), by the operator arrow
(\Cref{sec:grg-operator-arrow}), and by the cone dynamics of Supplement B.

Membership in the family is a statement in a chosen chart, and nothing here
selects a chart other than the link-trivializing one of
\Cref{hyp:gauss-flat-comparison}. The family is stated for the belief sector
at fixed model latents and not for the joint belief--model precision, and no
conditioning bound for the family is asserted. \status{NOT-CLAIMED}

\section{The interaction family}
\label{sec:gauss-interaction-family}

Throughout, agents are \(i\in V=\{1,\dots,N\}\), belief latents are
\(k_i\in\R^K\), and the stacked latent has dimension \(n=NK\).

\definitionheading{Interaction family}{def:gauss-interaction-family}
A matrix \(\Lambda\in\Sym^{NK}\), partitioned into \(K\times K\) agent blocks,
belongs to the interaction family \(\mathcal I(V)\) when there are self terms
\(A_i\in\PSD^{K}\) and edge weights \(W_{ij}=W_{ji}\in\PSD^{K}\), \(i\neq j\),
with
\begin{equation}
\Lambda_{ii}=A_i+\sum_{j\neq i}W_{ij},
\qquad
\Lambda_{ij}=-W_{ij}\quad(i\neq j).
\label{eq:gauss-interaction-family}
\end{equation}
Its Laplacian part is \(L:=\Lambda-\operatorname{blockdiag}(A_1,\dots,A_N)\).

The edge weights are matrices, so \(L\) is a matrix-weighted graph Laplacian,
the object of the matrix-weighted consensus literature
\citep{Trinh2018MatrixWeighted}.

\propositionheading{Energy and kernel}{prop:gauss-interaction-energy-kernel}
For \(\Lambda\in\mathcal I(V)\) and \(y=(y_1,\dots,y_N)\),
\begin{equation}
y^\top\Lambda y
=\sum_{i\in V}y_i^\top A_iy_i
+\sum_{i<j}\left(y_i-y_j\right)^\top W_{ij}\left(y_i-y_j\right).
\label{eq:gauss-interaction-energy}
\end{equation}
Consequently \(\Lambda\succeq0\), the consensus subspace
\(\mathbf 1\otimes\R^{K}\) lies in \(\ker L\), and \(\Lambda\succ0\) if and only
if the only \(y\) with \(A_iy_i=0\) for all \(i\) and
\(W_{ij}(y_i-y_j)=0\) for all \(i<j\) is \(y=0\). A necessary condition for
\(\Lambda\succ0\) is \(\sum_iA_i\succ0\). \status{ESTABLISHED}

\noindent\emph{Proof.}
Each unordered pair contributes
\(y_i^\top W_{ij}y_i+y_j^\top W_{ij}y_j-2y_i^\top W_{ij}y_j\), placing \(W_{ij}\)
on the two diagonal blocks and \(-W_{ij}\) on the two off-diagonal blocks;
adding the self terms reproduces \eqref{eq:gauss-interaction-family}. Every
term is nonnegative, and for \(M\succeq0\), \(v^\top Mv=0\) exactly when
\(Mv=0\), which gives the kernel criterion. For \(y=\mathbf 1\otimes v\) every
difference vanishes and the energy is \(v^\top(\sum_iA_i)v\). \(\square\)

With scalar weights a connected graph forces \(\ker L\) to be the consensus
subspace. A singular matrix weight leaves \(y_i-y_j\) free along
\(\ker W_{ij}\), so \(\ker L\) can be larger on a connected graph, and the
criterion of \Cref{prop:gauss-interaction-energy-kernel} replaces
connectivity.

\propositionheading{The interaction family is thin}{prop:gauss-interaction-family-thin}
For \(K\geq2\) and \(N\geq2\), \(\mathcal I(V)\) lies in a linear subspace of
\(\Sym^{NK}\) of codimension at least \(\binom N2\tfrac{K(K-1)}{2}\), hence has
Lebesgue measure zero there. Membership is not invariant under block-diagonal
changes of coordinates: the congruence of the member with \(N=K=2\),
\(A_1=A_2=I_2\), \(W_{12}=I_2\) by
\(\mathbf P=\operatorname{diag}(I_2,S)\),
\(S=\begin{psmallmatrix}1&1\\0&1\end{psmallmatrix}\), has off-diagonal block
\(-S\), which is not symmetric. \status{ESTABLISHED}

\noindent\emph{Proof.}
Membership requires each off-diagonal block \(\Lambda_{ij}\) to be symmetric as
a \(K\times K\) matrix, since \(W_{ij}=W_{ji}\) and \(W_{ji}=W_{ij}^\top\); these
are \(\tfrac{K(K-1)}2\) independent linear equations per unordered pair, and the
conic conditions only shrink the set. For the witness,
\((\mathbf P^\top\Lambda \mathbf P)_{12}=I_2(-I_2)S=-S\). \(\square\)

\section{When the interaction form holds}
\label{sec:gauss-hypothesis}

\subsection{Trivialized coordinates and flat comparison}

Write \(\Theta_{ij}:=\rho_b(\Theta_{ij}^b)\in\Gamma_b\subseteq\GL(K)\) for
the represented belief link of the edge copy \(j\to i\), the source copy of
row \(i\) when attention rows are present (\Cref{def:geo-graph-links}), with
\(\Gamma_b\) of \Cref{hyp:gen-lg-frame-group}. The residual between two agents is
\(k_i-\Theta_{ij}k_j\), not \(k_i-k_j\), so the plain differences of
\eqref{eq:gauss-interaction-energy} are a statement in coordinates in which
the links have been removed. \Cref{prop:gauss-coboundary-trivialization}
decides when that is possible: frames with \(U_i^{-1}\Theta_{ij}U_j=I\) exist
exactly when every cycle product of the links is the identity, and they are
then unique up to one common right factor on each connected component.

\hypothesisheading{Flat comparison}{hyp:gauss-flat-comparison}
The represented belief links form a graph coboundary: there are frames
\(U_i\in\Gamma_b\) with \(\rho_b(\Theta^b_e)=U_iU_j^{-1}\) on every edge copy
\(e:j\to i\), so the two copies of a mutually attending pair carry
\(\Theta_{ji}=\Theta_{ij}^{-1}\). In the coordinates \(z_i=U_i^{-1}k_i\), applied equally to latents and to means,
every residual is a plain difference,
\begin{equation}
k_i-\Theta_{ij}k_j=U_i\left(z_i-z_j\right).
\label{eq:gauss-flat-trivialization}
\end{equation}
The hypothesis excludes exactly the link assignments with a nontrivial cycle
product, including the two-cycle product \(\Theta_{ij}\Theta_{ji}\) of a
mutually attending pair, which the graph links of \Cref{def:geo-graph-links}
permit. It states
\Cref{def:gauss-interaction-family} in one chart and is the sufficient domain
of \Cref{prop:cg-gaussian-aggregation}.

\lemmaheading{Internal-edge annihilation}{lem:gauss-internal-edge}
Let \(\check\Theta_{ij}\) be the represented link in a chart common to \(i\)
and \(j\), and let \(W_{ij}\succeq0\). Identifying \(z_i=z_j=z\) leaves the edge
energy \(z^\top(I-\check\Theta_{ij})^\top W_{ij}(I-\check\Theta_{ij})z\),
which vanishes for every \(z\) if and only if
\(\operatorname{ran}(I-\check\Theta_{ij})\subseteq\ker W_{ij}\). For
\(W_{ij}\succ0\) this means \(\check\Theta_{ij}=I\); a singular weight can
hide a nontrivial link in its kernel. \status{ESTABLISHED}

\noindent\emph{Proof.}
The energy is \(\lVert W_{ij}^{1/2}(I-\check\Theta_{ij})z\rVert^2\), and
\(\ker W_{ij}^{1/2}=\ker W_{ij}\). \(\square\)

\remarkheading{Membership depends on the frame}{rem:gauss-membership-frame-dependence}
The change of variables \(y=Uz\), \(U=\operatorname{blockdiag}(U_1,\dots,U_N)\),
acts on precisions by the congruence \(\Lambda^{(z)}=U^\top\Lambda^{(y)}U\),
which preserves inertia by Sylvester's law \citep[Ch.~4]{HornJohnson2013} but
not membership in \(\mathcal I(V)\), since the block
\(U_i^\top\Lambda^{(y)}_{ij}U_j\) need not be symmetric when \(U_i\neq U_j\)
(\Cref{prop:gauss-interaction-family-thin}). A frame-independent criterion is
therefore invariant under block reframings with blocks in \(\Gamma_b\). Two
families of such invariants exist: generalized eigenvalues of the pair
\((L,\Lambda)\), which transform by one congruence
(\Cref{prop:ig-generalized-spectrum-invariance}), and law functionals of
unions of agent blocks, such as determinant gaps and mutual informations,
which a reframing leaves unchanged because it changes no law
(\Cref{sec:ig-gauge}). Eigenvalues of either matrix alone are invariant when
\(\Gamma_b\subseteq\operatorname{O}(K)\) (\Cref{prop:gen-lg-compact-frames}) and
frame dependent under \(\Gamma_b=\GL^+(K)\)
(\Cref{prop:ig-frame-dependent-spectra-determinants}). \status{ESTABLISHED}

\subsection{What a directed linear-Gaussian model contributes}

For the belief sector \eqref{eq:gauss-receiving-kernel} with transitions
\(\Omega_i=F_i\Theta_{i\operatorname{par}(i)}\), the blocks
\eqref{eq:gauss-directed-blocks} show that membership requires the candidate
weight
\begin{equation}
W_{i\operatorname{par}(i)}:=\Lambda_i\Omega_i=\Lambda_iF_i\Theta_{i\operatorname{par}(i)}
\label{eq:gauss-candidate-weight}
\end{equation}
to be symmetric and positive semidefinite. Under
\Cref{hyp:gauss-flat-comparison}, with \(j=\operatorname{par}(i)\), set
\begin{equation}
\widetilde\Omega_i:=U_i^{-1}\Omega_iU_j=U_i^{-1}F_iU_i,
\qquad
\widetilde\Lambda_i:=U_i^\top\Lambda_iU_i,
\label{eq:gauss-trivialized-transition}
\end{equation}
the second form of \(\widetilde\Omega_i\) using \(\Theta_{ij}=U_iU_j^{-1}\). The
links are then removed, and any remaining obstruction is attributable to the
gain.

\propositionheading{Two global-condition witnesses in the link-trivializing chart}{prop:gauss-global-condition-witnesses}
Under \Cref{hyp:gauss-flat-comparison}, the assembled precision in the chart
\(z=U^{-1}k\) is \(U^\top\Lambda U\), whose blocks are
\eqref{eq:gauss-directed-blocks} with \((\Lambda_i,\Omega_i)\) replaced by
\((\widetilde\Lambda_i,\widetilde\Omega_i)\). The chart is unique up to
\(U_i\mapsto U_ig\), \(g\in\Gamma_b\), independently on each connected component, which acts on its blocks by
\(M\mapsto g^\top Mg\), so symmetry of an off-diagonal block and the sign of a
symmetric block are the same in every such chart. In this chart there are
admissible parameters outside \(\mathcal I(V)\):
\begin{equation}
\widetilde\Lambda_i=\begin{psmallmatrix}2&1\\1&3\end{psmallmatrix}\succ0,
\qquad
\widetilde\Omega_i=\begin{psmallmatrix}0&1\\-1&0\end{psmallmatrix}\in\GL^{+}(2)
\qquad\Longrightarrow\qquad
-\widetilde\Lambda_i\widetilde\Omega_i=\begin{psmallmatrix}1&-2\\3&-1\end{psmallmatrix},
\label{eq:gauss-asymmetry-witness}
\end{equation}
which differs from its transpose by a matrix of Frobenius norm \(5\sqrt2\), so
no symmetric \(W\) exists; and \(K=1\), \(F_i=-1\), \(\Theta_{ij}=\Lambda_i=1\)
gives the symmetric block with \(W=-1<0\), as does \(F_i=-I\),
\(\Theta_{ij}=\Lambda_i=I\) with \(W=-I\) at every \(K\). \status{ESTABLISHED}

\noindent\emph{Proof.}
The residual of each factor is \(U_i(z_i-\widetilde\Omega_iz_j)\), so the
energy in \(z\) is that of \eqref{eq:gauss-directed-blocks} with the tilded
parameters. Uniqueness up to a common right factor on each connected component is part of
\Cref{prop:gauss-coboundary-trivialization}; under it
\(\widetilde\Omega_i\mapsto g^{-1}\widetilde\Omega_ig\) and
\(\widetilde\Lambda_i\mapsto g^\top\widetilde\Lambda_ig\), so
\(-\widetilde\Lambda_i\widetilde\Omega_i\mapsto g^\top(-\widetilde\Lambda_i\widetilde\Omega_i)g\);
a matrix \(M\) is symmetric exactly when \(g^\top(M-M^\top)g=0\), and a
symmetric \(M\) satisfies \(M\preceq0\) exactly when \(g^\top Mg\preceq0\). The
witnesses are direct computation: \(\det\widetilde\Lambda_i=5\),
\(\det\widetilde\Omega_i=1\), and the antisymmetric part of the displayed
block is \(\begin{psmallmatrix}0&-5\\5&0\end{psmallmatrix}\). \(\square\)

\propositionheading{Gain absorption and the orientation criterion}{prop:gauss-gain-absorption}
(o) The directed law depends on gain and link only through \(\Omega_i\), and a
block reframing \(k_i=\mathbf P_iz_i\) sends \(\Omega_i\) to
\(\mathbf P_i^{-1}\Omega_i\mathbf P_{\operatorname{par}(i)}\), the transformation law of a link.
Assume \Cref{hyp:gen-lg-frame-group}, so that \(\mathbf P_i\in\GL^+(K)\).
(i) On a directed forest whose transitions are invertible with
\(\det\Omega_i>0\), the blocks \(\mathbf P_r=I\) at every root and
\(\mathbf P_i=\Omega_i\mathbf P_{\operatorname{par}(i)}\) otherwise lie in \(\GL^+(K)\), and in
the coordinates \(z_i=\mathbf P_i^{-1}k_i\) the assembled precision lies in
\(\mathcal I(V)\) with
\begin{equation}
W_{i\operatorname{par}(i)}=\mathbf P_i^\top\Lambda_i\mathbf P_i,
\qquad
A_r=\Lambda_r\quad\text{at roots},
\qquad
A_i=0\quad\text{otherwise}.
\label{eq:gauss-gain-absorbed}
\end{equation}
Observation terms \(\operatorname{blockdiag}(O_1,\dots,O_N)\) with
\(O_i\succeq0\) keep it in \(\mathcal I(V)\), with self terms
\(A_i+\mathbf P_i^\top O_i\mathbf P_i\).
(ii) Let \(\Lambda\in\Sym_{++}^{2K}\) be a two-agent precision with invertible
off-diagonal block \(\Lambda_{12}\). There are \(\mathbf P_1,\mathbf P_2\in\GL^+(K)\) with
\(\operatorname{diag}(\mathbf P_1,\mathbf P_2)^\top\Lambda\operatorname{diag}(\mathbf P_1,\mathbf P_2)\in\mathcal I(\{1,2\})\)
if and only if
\begin{equation}
(-1)^K\det\Lambda_{12}>0.
\label{eq:gauss-orientation-criterion}
\end{equation}
(iii) A directed forest with invertible transitions has a chart reached by
frames in \(\GL^+(K)\) in which its assembled precision lies in
\(\mathcal I(V)\) if and only if \(\det\Omega_i>0\) for every nonroot \(i\).
(iv) If orientation-reversing frames are admitted, \(\mathbf P_i\in\GL(K)\), every
two-agent precision with invertible \(\Lambda_{12}\) is reached.
(v) The witness \eqref{eq:gauss-asymmetry-witness} has \(\det\widetilde\Omega_i=1\)
and is absorbed: with \(j\) its root, the blocks \(\mathbf P_i=I\),
\(\mathbf P_j=\widetilde\Omega_i^{-1}\) give the pure agreement factor with
\(W=\widetilde\Lambda_i\), \(A_i=0\), and
\(A_j=\widetilde\Omega_i^{-\top}\widetilde\Lambda_j\widetilde\Omega_i^{-1}\) for
the root precision \(\widetilde\Lambda_j\). For even \(K\) the gain \(-I\) has
determinant one and is absorbed in the same way. \status{ESTABLISHED}

\noindent\emph{Proof.}
(o) In the new coordinates the residual is
\[
\mathbf P_iz_i-\Omega_i\mathbf P_{\operatorname{par}(i)}z_{\operatorname{par}(i)}
=\mathbf P_i(z_i-\mathbf P_i^{-1}\Omega_i\mathbf P_{\operatorname{par}(i)}z_{\operatorname{par}(i)}).
\]
(i) Each \(\mathbf P_i\) is a product of matrices of positive determinant, and the
residual becomes \(\mathbf P_i(z_i-z_{\operatorname{par}(i)})\), so the energy is
\(\tfrac12\sum_i(z_i-z_{\operatorname{par}(i)})^\top \mathbf P_i^\top\Lambda_i\mathbf P_i(z_i-z_{\operatorname{par}(i)})
+\tfrac12\sum_rz_r^\top\Lambda_rz_r\), which is
\eqref{eq:gauss-interaction-energy} with the parameters
\eqref{eq:gauss-gain-absorbed}; an observation term adds
\(z_i^\top \mathbf P_i^\top O_i\mathbf P_iz_i\). (ii) If the reframed off-diagonal block is
\(-W\) with \(W\succeq0\) symmetric, invertibility of \(\Lambda_{12}\) makes
\(W\succ0\), and \(\det \mathbf P_1\det\Lambda_{12}\det \mathbf P_2=(-1)^K\det W\) forces
\eqref{eq:gauss-orientation-criterion}. Conversely, put \(\mathbf P_1=I\) and
\(\mathbf P_2=-\Lambda_{12}^{-1}\Lambda_{11}\), whose determinant
\((-1)^K\det\Lambda_{11}/\det\Lambda_{12}\) is positive under
\eqref{eq:gauss-orientation-criterion}. The reframed off-diagonal block is
\(-\Lambda_{11}\), so \(W=\Lambda_{11}\succ0\) and \(A_1=0\), while
\(A_2=\Lambda_{11}\Lambda_{12}^{-\top}\Lambda_{22}\Lambda_{12}^{-1}\Lambda_{11}-\Lambda_{11}\)
is the congruence by \(\Lambda_{12}^{-1}\Lambda_{11}\) of the Schur complement
\(\Lambda_{22}-\Lambda_{12}^\top\Lambda_{11}^{-1}\Lambda_{12}\succ0\).
(iii) The block of the edge to the parent is \(-\Lambda_i\Omega_i\), and
\((-1)^K\det(-\Lambda_i\Omega_i)=\det\Lambda_i\det\Omega_i\); a block reframing
multiplies that determinant by \(\det \mathbf P_i\det \mathbf P_{\operatorname{par}(i)}>0\), so
the necessity argument of (ii) applies edge by edge, and (i) gives sufficiency.
(iv) The construction of (ii) with \(\mathbf P_2=-\Lambda_{12}^{-1}\Lambda_{11}\) uses no
sign condition once \(\mathbf P_2\in\GL(K)\) is allowed. (v) With
\(z_j=\widetilde\Omega_ik_j\) the residual is \(z_i-z_j\) and the root prior
becomes \(\tfrac12z_j^\top\widetilde\Omega_i^{-\top}\widetilde\Lambda_j\widetilde\Omega_i^{-1}z_j\);
these blocks differ from those of (i) by the common right factor
\(\widetilde\Omega_i^{-1}\). \(\square\)

\propositionheading{Chart-independent witnesses}{prop:gauss-law-level-witnesses}
(i) Under \Cref{hyp:gen-lg-frame-group}, a single receiving factor with
invertible transition of negative determinant, for instance
\(\Omega_i=\operatorname{diag}(-1,1,\dots,1)\), has an assembled precision
outside \(\mathcal I(\{i,j\})\) in every chart reached by block reframings in
\(\GL^+(K)\).
(ii) For every cycle \(i_0,i_1,\dots,i_{n_\circ}=i_0\) of \(n_\circ\geq3\)
distinct agents whose blocks \(\Lambda_{i_ti_{t+1}}\) are invertible, the sign
of \(\prod_{t}\det\Lambda_{i_ti_{t+1}}\) is invariant under every block
reframing with \(\mathbf P_i\in\GL(K)\), and membership in \(\mathcal I(V)\) forces it
to equal \((-1)^{Kn_\circ}\).
(iii) Let a child \(i\) have two parents \(j,j'\) with invertible transitions
\(\Omega_j,\Omega_{j'}\), receiving precision \(\Lambda_i\), and proper root
priors, the directed acyclic generalization of \eqref{eq:gauss-receiving-kernel}
in which parent terms add. For odd \(K\) its assembled precision lies outside
\(\mathcal I(\{i,j,j'\})\) in every chart reached by block reframings in
\(\GL(K)\), orientation reversing or not. \status{ESTABLISHED}

\noindent\emph{Proof.}
(i) By \eqref{eq:gauss-orientation-criterion}, since
\((-1)^K\det(-\Lambda_i\Omega_i)=\det\Lambda_i\det\Omega_i<0\). (ii) A block
reframing multiplies \(\det\Lambda_{i_ti_{t+1}}\) by
\(\det \mathbf P_{i_t}\det \mathbf P_{i_{t+1}}\), and each agent of the cycle occurs in two
factors, so the product is multiplied by \(\prod_t(\det \mathbf P_{i_t})^2>0\). In
\(\mathcal I(V)\) an invertible off-diagonal block is \(-W_e\) with
\(W_e\succ0\), so the product equals \((-1)^{Kn_\circ}\prod_e\det W_e\).
(iii) The factor
\(\tfrac12(k_i-\Omega_jk_j-\Omega_{j'}k_{j'})^\top\Lambda_i(k_i-\Omega_jk_j-\Omega_{j'}k_{j'})\)
gives \(\Lambda_{ij}=-\Lambda_i\Omega_j\), \(\Lambda_{ij'}=-\Lambda_i\Omega_{j'}\),
and \(\Lambda_{jj'}=\Omega_j^\top\Lambda_i\Omega_{j'}\), all invertible, while the
root priors add only to diagonal blocks. The product of the three determinants
around the triangle is
\((\det\Lambda_i)^3(\det\Omega_j)^2(\det\Omega_{j'})^2>0\), whereas (ii) requires
the sign \((-1)^{3K}=-1\) for odd \(K\). \(\square\)

The gain \(F_i\) is a free generative parameter, not required to be symmetric,
positive, or close to the identity. On forests the witnesses of
\Cref{prop:gauss-global-condition-witnesses} are properties of the
link-trivializing chart, and the law-level obstructions are an
orientation-reversing gain and, for odd \(K\), a v-structure
(\Cref{prop:gauss-law-level-witnesses}). The interaction family is therefore
not a consequence of the general directed model, and
\Cref{hyp:gauss-global-interaction} imposes it in the link-trivializing
coordinates.

\subsection{Exactly when the form holds}

\propositionheading{Exact global interaction condition}{prop:gauss-global-interaction-condition}
(i) The assembled prior precision \eqref{eq:gauss-directed-blocks} lies in
\(\mathcal I(V)\), necessarily with the edge weights
\eqref{eq:gauss-candidate-weight}, if and only if, for every nonroot \(i\),
\begin{equation}
W_i:=\Lambda_i\Omega_i\in\PSD^{K},
\label{eq:gauss-global-cond-1}
\end{equation}
and, for every agent \(i\),
\begin{equation}
A_i:=\Lambda_i-W_i+\sum_{i'\in\operatorname{ch}(i)}\left(W_{i'}\Lambda_{i'}^{-1}W_{i'}-W_{i'}\right)\succeq0,
\label{eq:gauss-global-cond-2}
\end{equation}
with the term \(-W_i\) omitted at a root.
(ii) With the observation kernel \eqref{eq:gen-lg-obs}, the belief conditional
\(p(k\given m,o,X)\) has precision
\begin{equation}
\Lambda^{\mathrm{post}}=\Lambda^{\mathrm{prior}}+\operatorname{blockdiag}(O_1,\dots,O_N),
\qquad
O_i:=\mathsf C_i^{b\top}(\Sigma_i^o)^{-1}\mathsf C_i^b\succeq0,
\label{eq:gauss-belief-conditional-precision}
\end{equation}
and \(\Lambda^{\mathrm{post}}\in\mathcal I(V)\) if and only if
\eqref{eq:gauss-global-cond-1} holds and \(A_i+O_i\succeq0\) at every agent,
while bridge terms, offsets, and the model part of each observation mean enter
only its information vector. The prior condition of (i) is therefore
sufficient for membership of the belief conditional.
(iii) It is not necessary: for one scalar receiving factor \(j\to i\) with
\(\Lambda_i=1\), \(\Omega_i=\tfrac12\), root precision \(\Lambda_j=\tfrac1{10}\),
and an observation on \(j\) with \(O_j=\tfrac12\), the prior precision
\(\begin{psmallmatrix}1&-1/2\\-1/2&7/20\end{psmallmatrix}\) is positive
definite, with determinant \(\tfrac1{10}\), and has
\(A_j=-\tfrac3{20}\), so it lies outside the family, while the belief
conditional has \(A_j+O_j=\tfrac7{20}\) and \(A_i=\tfrac12\) and lies inside it.
(iv) Factorwise membership is sufficient but not necessary for global
membership: for \(\Lambda_i=1\), \(\Omega_i=\tfrac12\), and \(\Lambda_j=1\), in
the block order \((i,j)\),
\begin{equation}
\Lambda
=
\begin{pmatrix}
1&-\tfrac12\\[2pt]
-\tfrac12&\tfrac54
\end{pmatrix}
=
\begin{pmatrix}
A_i+W&-W\\[2pt]
-W&A_j+W
\end{pmatrix},
\qquad
W=\tfrac12,\quad A_i=\tfrac12,\quad A_j=\tfrac34,
\label{eq:gauss-global-cancellation-witness}
\end{equation}
is positive definite and in \(\mathcal I(\{i,j\})\), although the receiving
factor alone is not: its parent-side residual is
\(\Omega_i^2\Lambda_i-W=-\tfrac14\), offset by the root prior.
\status{ESTABLISHED}

\noindent\emph{Proof.}
(i) Condition \eqref{eq:gauss-global-cond-1} is exactly that each off-diagonal
block has the form \(-W\) with \(W\) symmetric positive semidefinite, and the
off-diagonal blocks of a member determine its edge weights. Given it,
\(\Omega_{i'}=\Lambda_{i'}^{-1}W_{i'}\) and
\(\Omega_{i'}^\top\Lambda_{i'}\Omega_{i'}=W_{i'}\Lambda_{i'}^{-1}W_{i'}\) by
symmetry of \(W_{i'}\). The self term is the residual
\(A_i=\Lambda_{ii}-\sum_{j\sim i}W_{ij}\) over the edge to the parent, of weight
\(W_i\), and the edges to the children, of weights \(W_{i'}\); substituting
\eqref{eq:gauss-directed-blocks} gives \eqref{eq:gauss-global-cond-2}.
(ii) At fixed \(m\) and \(o\), the \(k\)-dependent terms of the negative log
joint density are the belief factors, in which \(m\) and the offsets enter only
the means, and the observation factors, each contributing
\(\tfrac12k_i^\top O_ik_i\) plus terms at most linear in \(k_i\). The Hessian in
\(k\) is therefore \eqref{eq:gauss-belief-conditional-precision}, with the
off-diagonal blocks of the prior and self-term residuals \(A_i+O_i\).
(iii) and (iv) are direct evaluation of \eqref{eq:gauss-global-cond-2}: in
(iii) \(A_j=\tfrac1{10}+\tfrac14-\tfrac12\) and \(A_i=1-\tfrac12\); in (iv) the
leading principal minors are \(1\) and \(1\), and \(A_j=1-\tfrac14\). \(\square\)

\propositionheading{Edge-local characterization with singular gains allowed}{prop:gauss-edge-local-characterization}
Consider a single receiving factor \(j\to i\) in isolation, contributing
\(\Lambda_i\) to the \(ii\) block, \(\Omega_i^\top\Lambda_i\Omega_i\) to the
\(jj\) block, and \(-\Lambda_i\Omega_i\) to the \(ij\) block. It lies in
\(\mathcal I(\{i,j\})\) if and only if
\begin{equation}
\Pi_i:=\Lambda_i^{-1/2}\left(\Lambda_i\Omega_i\right)\Lambda_i^{-1/2}
\text{ is an orthogonal projection,}
\label{eq:gauss-edge-local}
\end{equation}
equivalently if and only if \(\Lambda_i\Omega_i\) is symmetric and
\(\Omega_i\) is idempotent, that is \(\Omega_i\) is a \(\Lambda_i\)-orthogonal
projection. Its self terms are then \(A_i=\Lambda_i^{1/2}(I-\Pi_i)\Lambda_i^{1/2}\)
and \(A_j=0\), and the full-rank case \(\Pi_i=I\) is \(\Omega_i=I\), the pure
agreement factor with \(W=\Lambda_i\). The gains satisfying the condition
form a Lebesgue-null subset of \(\R^{K\times K}\). \status{ESTABLISHED}

\noindent\emph{Proof.}
Membership requires \(W:=\Lambda_i\Omega_i\) symmetric and positive
semidefinite, \(\Lambda_i-W\succeq0\), and \(\Omega_i^\top\Lambda_i\Omega_i-W\succeq0\).
With \(W\) symmetric, \(\Pi_i\) is symmetric; the first two conditions read
\(0\preceq \Pi_i\preceq I\), and since
\(\Omega_i^\top\Lambda_i\Omega_i=W\Lambda_i^{-1}W=\Lambda_i^{1/2}\Pi_i^2\Lambda_i^{1/2}\),
the third reads \(\Pi_i^2\succeq \Pi_i\). As \(0\preceq \Pi_i\preceq I\) implies
\(\Pi_i^2\preceq \Pi_i\), the conditions hold exactly when \(\Pi_i^2=\Pi_i\). Conversely every
orthogonal projection satisfies them, with \(\Omega_i=\Lambda_i^{-1/2}\Pi_i\Lambda_i^{1/2}\)
idempotent; the self terms follow by substitution. Idempotents of rank \(r\)
form the conjugation orbit of \(\operatorname{diag}(I_r,0)\), whose stabilizer
is \(\GL(r)\times\GL(K-r)\), a smooth submanifold of dimension
\(K^2-r^2-(K-r)^2=2r(K-r)<K^2\); the finite union over \(r\) is null.
\(\square\)

\corollaryheading{Invertible-gain factor}{cor:gauss-invertible-gain-factor}
In the link-trivializing chart, an isolated factor with invertible transition
lies in \(\mathcal I(\{i,j\})\) if and only if \(\Omega_i=I\). When every chart
reached by frames in \(\GL^+(K)\) is allowed, an invertible factor with
\(\det\Omega_i>0\) is instead absorbed by \Cref{prop:gauss-gain-absorption}.
\status{ESTABLISHED}

\noindent\emph{Proof.}
\Cref{prop:gauss-edge-local-characterization} makes \(\Omega_i\) idempotent,
and multiplying \(\Omega_i^2=\Omega_i\) by \(\Omega_i^{-1}\) gives \(\Omega_i=I\);
the converse is its full-rank case. \(\square\)

With singular gains a proper projection constrains only the components of the
residual inside \(\operatorname{range}(\Pi_i)\) and leaves the complement to the
receiving-side self term.

\hypothesisheading{Global interaction form}{hyp:gauss-global-interaction}
Where a result names this hypothesis, the precision it acts on belongs to
\(\mathcal I(V)\) in the link-trivializing coordinates. For the directed model
\eqref{eq:gauss-receiving-kernel} the precision is that of the belief
conditional \eqref{eq:gauss-belief-conditional-precision}, and by
\Cref{prop:gauss-global-interaction-condition}(ii) the hypothesis holds exactly
when \eqref{eq:gauss-global-cond-1} holds and \(A_i+O_i\succeq0\) at every
agent; the prior balance \eqref{eq:gauss-global-cond-2} is sufficient but not
necessary, and the receiving factors need not satisfy
\Cref{prop:gauss-edge-local-characterization}
(\eqref{eq:gauss-global-cancellation-witness}). The hypothesis licenses the
aggregation closure of \Cref{sec:cg-gaussian-aggregation} and the operator
arrow of \Cref{sec:grg-operator-arrow}, which are stated for the family, on a
precision assembled from the directed model. It excludes precisions whose
off-diagonal blocks are asymmetric or sign-incorrect in the link-trivializing
chart, which the directed model does not prevent
(\Cref{prop:gauss-global-condition-witnesses}) and which form all but a
null set of precisions when \(K\geq2\) and \(N\geq2\)
(\Cref{prop:gauss-interaction-family-thin}).

\remarkheading{Factorwise-agreement realization}{hyp:gauss-factorwise-agreement}
Requiring every receiving factor to satisfy
\Cref{prop:gauss-edge-local-characterization}, so that the interaction
parameters are read from the generative model one factor at a time, is
strictly stronger than \Cref{hyp:gauss-global-interaction}: it admits
\(\Lambda_i\)-orthogonal projections with singular gains and only
\(\Omega_i=I\) with invertible gains, and it excludes the global member
\eqref{eq:gauss-global-cancellation-witness}. \status{ESTABLISHED}

With latent links in the exponential chart, the matter precision of
\Cref{sec:glink-induced} is a function \(\Lambda(\mathring\omega)\) of the
generators that belongs to this family at zero generators, and at zero
information vector one half its log-determinant is, up to a constant, the
induced action of \Cref{prop:glink-induced-action}.

\section{Nonemptiness}
\label{sec:gauss-nonempty}

\propositionheading{Constructive nonemptiness and the kernel}{prop:gauss-interaction-nonempty}
Under \Cref{hyp:gauss-flat-comparison}, let each edge of a graph \((V,E)\)
carry an agreement factor with residual \(\mu_i-\Theta_{ij}\mu_j\) and link
covariance \(\Sigma^{\mathrm{link}}_{ij}\in\Sym_{++}^{K}\), and let each vertex
carry a quadratic self term \(A_i\succeq0\) in the trivialized coordinates; a
self term is the precision of a proper full-dimensional prior only when
\(A_i\succ0\). The assembled trivialized energy is
\begin{equation}
\tfrac12z^\top\Lambda z=\tfrac12\sum_{i}z_i^\top A_iz_i
+\tfrac12\sum_{(i,j)\in E}\left(z_i-z_j\right)^\top W_{ij}\left(z_i-z_j\right),
\qquad
W_{ij}=U_i^\top(\Sigma^{\mathrm{link}}_{ij})^{-1}U_i,
\label{eq:gauss-nonempty-construction}
\end{equation}
so \(\Lambda\in\mathcal I(V)\cap\PSD^{NK}\), every edge weight is positive
definite, and
\begin{equation}
\ker\Lambda
=\left\{z:A_iz_i=0\text{ for every }i,
\quad z_i=z_j\text{ for every }(i,j)\in E\right\}.
\label{eq:gauss-nonempty-kernel}
\end{equation}
If \(V_\alpha\) are the connected components of \((V,E)\), equivalently
\begin{equation}
\ker\Lambda
=\bigoplus_\alpha
\left\{\mathbf1_{V_\alpha}\otimes v:
v\in\bigcap_{i\in V_\alpha}\ker A_i\right\}.
\label{eq:gauss-nonempty-component-kernel}
\end{equation}
Consequently \(\Lambda\succ0\) exactly when
\(\bigcap_{i\in V_\alpha}\ker A_i=\{0\}\) for every component; one
positive-definite self term per component suffices, and taking every
\(A_i\succ0\) constructs an element of \(\mathcal I(V)\cap\Sym_{++}^{NK}\) whose
self terms are proper priors. \status{ESTABLISHED}

\noindent\emph{Proof.}
By \eqref{eq:gauss-flat-trivialization} the residual is \(U_i(z_i-z_j)\), so the
edge quadratic form has the matrix
\(U_i^\top(\Sigma^{\mathrm{link}}_{ij})^{-1}U_i\) between two copies of
\(z_i-z_j\), and that matrix is symmetric positive definite because \(U_i\) is
invertible.
Comparison with \eqref{eq:gauss-interaction-energy} identifies the parameters.
The energy vanishes exactly when \(A_iz_i=0\) at every vertex and, since every
\(W_{ij}\succ0\), \(z_i=z_j\) on every edge; the latter makes \(z\) constant on
each component. A positive-semidefinite matrix is positive definite exactly
when its kernel is trivial. \(\square\)

Without trivialization the same agreement factor contributes the off-diagonal
block \(-(\Sigma^{\mathrm{link}}_{ij})^{-1}\Theta_{ij}\), which is symmetric only
in special position; the trivialization is what places the construction in
the family.
